\chapter{DATA ANALYSIS -- POST-COVID VIEWERSHIP TRANSFORMATION}
\label{ch:viewership}

\section{Introduction}
\label{sec:viewer_intro}

This chapter presents the analysis of the 400-respondent
viewer-survey data collected between October 2024 and February 2025,
following the design described in Chapter~\ref{ch:methodology}.  The
chapter is organised around the 17 substantive variables of the
viewer questionnaire and is structured to (a) describe the sample,
(b) compare pre- and post-COVID values of each behavioural variable
through paired tests of statistical significance, (c) examine
demographic moderators where present, (d) interpret the results
against the literature reviewed in Chapter~\ref{ch:literature}, and
(e) summarise the standing of hypotheses H1--H5.

All analyses were performed in IBM SPSS Statistics 28; selected
multiple-regression models were replicated in R 4.4.0
\citep{rcoreteam2024}.  Effect sizes are reported using Cohen's $d$
for paired-sample $t$-tests and partial $\eta^2$ for ANOVA.
Statistical significance is reported at conventional thresholds.
The complete numeric tables underpinning the figures are
reproduced in Appendix~C.

\section{Profile of the Viewer Sample}
\label{sec:profile}

\subsection{Age, gender and location}

The age distribution of the 400 respondents (Figure~\ref{fig:age})
reflects the Tamil-Nadu adult population skew toward the 18--34 age
band, which together account for 62\% of the sample.  The gender
composition (Figure~\ref{fig:gender}) is 55\% male, 44\% female and
1\% other / no-response, aligned with our quota plan and consistent
with Ormax-reported gender distribution of Tamil cinema-goers
\citep{ormax2023}.  The geographic distribution
(Figure~\ref{fig:location}) places the bulk of the sample in
Chennai (104) and the seven major Tier-1 / Tier-2 cities of Tamil
Nadu, with smaller cells for Tier-2 / Tier-3 (68) and rural (40).

\begin{figure}[h]
\centering
\includegraphics[width=0.78\textwidth]{fig_01_age.png}
\caption{Age distribution of viewer respondents (n=400).}
\label{fig:age}
\end{figure}

\begin{figure}[h]
\centering
\includegraphics[width=0.55\textwidth]{fig_02_gender.png}
\caption{Gender composition.}
\label{fig:gender}
\end{figure}

\begin{figure}[h]
\centering
\includegraphics[width=0.85\textwidth]{fig_03_location.png}
\caption{Geographic distribution of respondents.}
\label{fig:location}
\end{figure}

\subsection{Income and education}

Household monthly income (Figure~\ref{fig:income}) is concentrated in
the INR 25{,}000--100{,}000 band (58\% combined), reflecting both the
sampling effort to span middle-income strata and the contemporary
income distribution of urban Tamil Nadu.  Educational qualification
(Figure~\ref{fig:education}) is dominated by undergraduate-degree
holders (46\%) and post-graduate qualification (27\%), reflecting the
relatively high literacy and tertiary-enrolment of the state
\citep{census2011}.

\begin{figure}[h]
\centering
\includegraphics[width=0.78\textwidth]{fig_04_income.png}
\caption{Monthly household income (INR).}
\label{fig:income}
\end{figure}

\begin{figure}[h]
\centering
\includegraphics[width=0.78\textwidth]{fig_05_education.png}
\caption{Educational qualification.}
\label{fig:education}
\end{figure}

\subsection{Sample-quality assessment}

The achieved demographic distribution closely mirrors the planned
quotas (Table~\ref{tab:quotas}, Chapter~\ref{ch:methodology}).  No
single demographic cell shows divergence from the quota target by
more than 2 percentage points, supporting the representativeness
claim within the defined population frame.  Missing data are minimal
(below 1.4\% on any single item) and were addressed by listwise
deletion in the relevant inferential models.

\section{Theatre Attendance: Pre- versus Post-COVID}
\label{sec:theatre_results}

Mean monthly theatre visits in 2019 ($M = 3.10$, $SD = 1.65$) fell
sharply to a 2024 mean of $M = 1.46$ ($SD = 1.21$), a difference of
$-1.64$ visits per month or a 53\% decline.  A paired-sample
$t$-test indicates the difference is highly statistically
significant: $t(399) = 26.3$, $p < 0.001$.  Cohen's $d = 1.31$ --
a \emph{large} effect size by conventional benchmarks
\citep{cohen1988}.

Disaggregation by age group (Figure~\ref{fig:theatre_age}) reveals
a non-uniform decline.  The 35--44 age band exhibits the steepest
absolute fall (3.4 \(\to\) 1.5 visits, $-1.9$), while the 55+ band
exhibits a smaller absolute but proportionally similar decline.
The 18--24 band, although still the most cinema-going age, also
declined by 1.5 visits per month.

\begin{figure}[h]
\centering
\includegraphics[width=0.85\textwidth]{fig_06_theatre_age.png}
\caption{Mean monthly theatre visits -- pre vs post-COVID, by age.}
\label{fig:theatre_age}
\end{figure}

A factorial ANOVA with age (5 levels) and location-tier (Chennai vs
non-Chennai) as factors and the pre-post difference as the dependent
variable found significant main effects of age, $F(4,390) = 4.26$,
$p = 0.002$, $\eta^2_p = 0.042$, and of location-tier, $F(1,390) =
8.01$, $p = 0.005$, $\eta^2_p = 0.020$, with no significant
interaction.  Post-hoc Tukey HSD comparisons indicate the 35--44 vs
18--24 contrast as the largest pairwise difference, $p = 0.011$.

\textbf{Implication.}  H1 is supported.  Theatrical attendance has
fallen significantly across all demographic strata.  The age
gradient is consistent with the global pattern reported by
\citet{schauerte2021} for European markets and by
\citet{garcia2023streaming} for North America.

\section{OTT Subscription Adoption}
\label{sec:ott_results}

The mean number of OTT subscriptions held by respondents rose from
$M = 1.12$ in 2019 ($SD = 0.93$) to $M = 2.81$ in 2024
($SD = 1.42$), a 151\% increase
(Figure~\ref{fig:ott_subs}).  The paired difference is highly
significant: $t(399) = 24.8$, $p < 0.001$, Cohen's $d = 1.24$.

\begin{figure}[h]
\centering
\includegraphics[width=0.85\textwidth]{fig_07_ott_subs.png}
\caption{OTT subscription holdings -- before vs after pandemic.}
\label{fig:ott_subs}
\end{figure}

McNemar's tests on the binary state ``holds 0 subscriptions'' versus
``holds 1+'' show a highly significant migration: 96 respondents who
held zero in 2019 reported 1+ in 2024, against only 7 in the
opposite direction ($\chi^2 = 76.5$, $p < 0.001$).

\textbf{Implication.}  H2 is supported.  OTT adoption has more than
doubled, with the share of zero-subscription households falling from
36\% to 4\%.  Diffusion in this Tamil-Nadu sample is approaching
late-majority saturation as defined by \citet{rogers2003}.

\section{OTT Platform Preference for Tamil Films}
\label{sec:platforms_results}

Respondents were invited to indicate (multi-response) which
platforms they regularly used for Tamil-language film viewing.
Figure~\ref{fig:platforms} reports the responses ranked from highest
to lowest.

\begin{figure}[h]
\centering
\includegraphics[width=0.92\textwidth]{fig_08_platforms.png}
\caption{Preferred OTT platforms for Tamil-language films
(multi-response, n=400).}
\label{fig:platforms}
\end{figure}

Three observations are noteworthy.  First, Amazon Prime Video (312
respondents, 78\%) and Sun NXT (268 respondents, 67\%) emerge as the
two most-used platforms, reflecting respectively the global SVOD's
strong Tamil content investment and Sun's deep regional catalogue.
Second, YouTube (free) is the third most-used platform (226, 56\%),
underscoring the persistence of free / AVOD consumption alongside
SVOD.  Third, Aha Tamil (142, 36\%) -- the dedicated regional SVOD
-- is below the median for our sample, although it has shown the
fastest growth from a near-zero pre-pandemic base \citep{
aha2022launchreport,ficciey2024}.

\section{Monthly Entertainment Spend}
\label{sec:spend_results}

Monthly household spend on entertainment rose from a 2019 mean of
INR 622 ($SD = 281$) to a 2024 mean of INR 736 ($SD = 354$), an 18\%
increase (Figure~\ref{fig:spend}).  The paired-sample test is
significant: $t(399) = 8.41$, $p < 0.001$, $d = 0.42$ -- a
\emph{small-to-medium} effect.  The mix shifted from 71\%
theatrical / 18\% subscription / 11\% other in 2019 to 38\%
theatrical / 51\% subscription / 11\% other in 2024.

\begin{figure}[h]
\centering
\includegraphics[width=0.78\textwidth]{fig_09_spend.png}
\caption{Monthly household entertainment spend -- pre vs post-COVID.}
\label{fig:spend}
\end{figure}

The implication for the industry's revenue mix is direct: the
absolute size of the household entertainment wallet has grown
modestly, but its allocation has decisively shifted toward
subscription, foreshadowing the supply-side findings of
Chapter~\ref{ch:industry_reconfig}.

\section{Weekly OTT Viewing Hours}
\label{sec:hours_results}

The mean number of hours per week dedicated to OTT viewing rose
from $M = 6.4$ (SD$=4.1$) in 2019 to $M = 12.2$ ($SD = 5.6$) in
2024, a 91\% increase (Figure~\ref{fig:ott_hours}).  Paired-sample
test: $t(399) = 27.6$, $p < 0.001$, $d = 1.38$.

\begin{figure}[h]
\centering
\includegraphics[width=0.85\textwidth]{fig_10_ott_hours.png}
\caption{Weekly OTT viewing hours.}
\label{fig:ott_hours}
\end{figure}

The post-pandemic distribution is also more right-skewed: a sub-population
of approximately 17\% of respondents now reports above 20 hours per
week, a category that was negligible in 2019.  This sub-population
corresponds approximately to the ``heavy SVOD user'' segment
identified by \citet{lotz2022,lim2020}.

\section{Genre Preference Shift}
\label{sec:genre_results}

Respondents allocated 100\% across ten genre categories representing
their typical Tamil-film viewing.  Figure~\ref{fig:genres} compares
the pre- and post-COVID distributions.

\begin{figure}[h]
\centering
\includegraphics[width=0.92\textwidth]{fig_11_genres.png}
\caption{Genre preference -- pre vs post-COVID.}
\label{fig:genres}
\end{figure}

Mass-action declined from 31\% to 24\%, family drama from 22\% to
16\%, and romance from 18\% to 13\%.  Conversely, thriller / crime
rose from 14\% to 22\%, realistic / indie from 9\% to 18\%,
investigative from 7\% to 14\% and anthology from 3\% to 11\%.

A chi-square test of homogeneity on the underlying frequency cross-tabulation
yields $\chi^2(9) = 89.4$, $p < 0.001$, indicating a statistically
significant shift in the genre-preference distribution.  Effect-size
estimation using Cramer's $V$ gives $V = 0.21$, a moderate
effect.

\textbf{Implication.}  H4 is supported.  The post-pandemic genre
distribution exhibits a clear migration from ``mass'' clusters to
``thriller / realistic / anthology'' clusters.  The finding mirrors
\citet{sundet2021}'s European data and \citet{lotz2022}'s
US/Anglophone analysis but with a sharper migration in our Tamil
sample, possibly reflecting catch-up to the global content-driven
norm.

\section{Co-Viewing Patterns}
\label{sec:co_viewing_results}

Co-viewing (Figure~\ref{fig:coviewing}) shifted markedly toward
solo consumption.  ``Alone'' rose from 22\% to 39\%, ``with family''
fell from 46\% to 31\%, and ``with friends'' fell from 20\% to 13\%.
``With partner'' rose modestly from 12\% to 17\%.

\begin{figure}[h]
\centering
\includegraphics[width=0.78\textwidth]{fig_12_coviewing.png}
\caption{Co-viewing context -- pre vs post-COVID.}
\label{fig:coviewing}
\end{figure}

A chi-square test of independence on the four-by-two cross-tabulation
yields $\chi^2(3) = 38.7$, $p < 0.001$.  The migration to solo
consumption is consistent with the device shift discussed in the
next section: solo consumption on smartphones has substituted for
some part of group / family consumption in earlier modes
(family-room television, theatre).

\section{Device of Consumption}
\label{sec:device_results}

Smartphone remains the dominant primary device but its share grew
modestly from 46\% to 52\% (Figure~\ref{fig:devices}).  The
significant change is the rise of smart television from 18\% to 27\%
and the corresponding decline of laptop / desktop (16\% \(\to\) 10\%)
and STB / cable (14\% \(\to\) 4\%).

\begin{figure}[h]
\centering
\includegraphics[width=0.85\textwidth]{fig_13_devices.png}
\caption{Primary device for Tamil-film viewing.}
\label{fig:devices}
\end{figure}

The implication is a household-level migration toward two
complementary devices: the personal smartphone for solo and
on-the-go consumption, and the smart television as the
\emph{social} viewing device that has replaced traditional
satellite-television family viewing.

\section{Theatrical Window Preference}
\label{sec:window_results}

Respondents were asked to indicate their preferred theatre-to-OTT
release-window length.  Figure~\ref{fig:window} reports the
distribution.

\begin{figure}[h]
\centering
\includegraphics[width=0.78\textwidth]{fig_14_window.png}
\caption{Preferred theatre-to-OTT release window (post-COVID).}
\label{fig:window}
\end{figure}

The modal preference is the 3--4 week window (31\%), followed by
5--8 weeks (24\%) and 1--2 weeks (22\%).  The pre-pandemic norm of
$>$8 weeks is preferred by only 12\% of respondents.  ``Same-day
OTT'' is preferred by 11\%.

This consumer-side finding is matched closely by the supply-side
distribution of \emph{actual} negotiated windows reported by
specialists in Chapter~\ref{ch:industry_reconfig}
(Figure~\ref{fig:window_negotiated}), evidencing convergence
between viewer expectation and producer / OTT practice.

\section{Drivers and Barriers of Theatre Return}
\label{sec:return_results}

For respondents who reported any theatre attendance in 2024, the
multi-response question on factors driving theatre return /
non-return produced the distribution in
Figure~\ref{fig:theatre_reasons}.

\begin{figure}[h]
\centering
\includegraphics[width=0.92\textwidth]{fig_15_theatre_reasons.png}
\caption{Drivers and barriers of theatre return.}
\label{fig:theatre_reasons}
\end{figure}

The strongest pull factor is the ``big-screen experience'' (62\%),
followed by ``star release / FDFS'' (44\%) and ``group / family
outing'' (39\%).  The strongest push factor away from theatre is
``convenience of OTT'' (51\%), followed by ``content not engaging''
(33\%) and ``lack of OTT release'' (28\%, indicating that some
viewers \emph{prefer} a film to be on OTT before they will engage
with it).  Notably, ``cleanliness / safety concerns'', a major early-pandemic
barrier, has fallen to 14\% by 2024, suggesting that COVID-specific
hygiene anxiety is no longer a structural barrier.

\section{Cross-Language Consumption}
\label{sec:language_results}

Figure~\ref{fig:languages} compares pre- and post-COVID cross-language
consumption.

\begin{figure}[h]
\centering
\includegraphics[width=0.92\textwidth]{fig_16_languages.png}
\caption{Cross-language film consumption.}
\label{fig:languages}
\end{figure}

The proportion of viewers reporting Tamil-only consumption fell from
42\% to 27\%.  The proportion reporting Tamil + English / foreign
content rose from 9\% to 28\% -- the largest single increase in any
language pair, driven by Netflix and Prime English-language
originals \citep{lotz2022,jenner2018}.  Tamil + Telugu and
Tamil + Malayalam also saw meaningful growth.

The finding is theoretically significant.  In a cultural-industries
framework \citep{hesmondhalgh2019,straubhaar2007cultural}, language
remains an important compatibility constraint.  The post-pandemic
softening of this constraint among Tamil viewers -- particularly the
younger demographic and the higher-education sub-sample -- suggests
that streaming has \emph{eroded} but not \emph{eliminated} the
language-compatibility friction.

\section{Time of Viewing}
\label{sec:time_results}

Time-of-day shifts (Figure~\ref{fig:time}) show a modest
re-distribution toward late-night (12--3 am) viewing, from 16\% to
24\%, and toward morning / afternoon time slots; the prime-time
9--12 pm slot fell from 47\% to 38\%.

\begin{figure}[h]
\centering
\includegraphics[width=0.85\textwidth]{fig_17_time.png}
\caption{Time-of-day of primary film viewing.}
\label{fig:time}
\end{figure}

The re-distribution is consistent with both work-from-home flexibility
and the asynchronous-consumption affordance of OTT.

\section{Star versus Content as Decision Drivers}
\label{sec:star_content_results}

Perhaps the single most theoretically charged finding of the survey
is presented in Figure~\ref{fig:star_content}.

\begin{figure}[h]
\centering
\includegraphics[width=0.62\textwidth]{fig_18_star_content.png}
\caption{Decision driver -- star vs content.}
\label{fig:star_content}
\end{figure}

In 2019, 64\% of respondents identified ``a film's lead star or
director'' as the primary driver of their decision to watch.  In
2024, this had fallen to 46\%.  Conversely, ``the film's
story / content / reviews'' rose from 36\% to 54\%.

McNemar's test on the binary star-driven / content-driven variable
yields $\chi^2 = 49.7$, $p < 0.001$.  The shift is statistically
robust.

\textbf{Implication.}  H3 is supported.  The post-pandemic Tamil
viewer is markedly less star-driven and more content-driven than
the pre-pandemic baseline.  The implication for star economics --
explored in Chapter~\ref{ch:industry_reconfig} via specialist
data on star-fee restructuring -- is direct: the demand-side
willingness to pay a star premium has weakened.  The finding does
not, however, imply that the star system has collapsed: the top-tier
``Vijay / Ajith / Kamal / Rajini-class'' star remains a major demand
pull, as the specialist data show.

\section{Self-Reported Piracy Consumption}
\label{sec:piracy_results}

Figure~\ref{fig:piracy} reports self-reported piracy frequency.  The
``Never'' category rose from 22\% to 38\% (a positive development
from the industry's perspective), and the ``Frequently'' category
fell from 19\% to 9\%.

\begin{figure}[h]
\centering
\includegraphics[width=0.78\textwidth]{fig_19_piracy.png}
\caption{Self-reported piracy consumption.}
\label{fig:piracy}
\end{figure}

The finding is consistent with the broader argument in
\citet{lobato2019,jin2019} that legitimate streaming alternatives
\emph{reduce} piracy when priced and curated appropriately.  The
caveat from Chapter~\ref{ch:methodology} on self-report bias
applies; the magnitude is plausibly under-reported, but the
\emph{direction} of change is well-supported by the trade-press
\citep{filmcompanion2022mid,behindwoods2023}.

\section{Willingness to Pay for TVOD Rentals}
\label{sec:tvod_results}

Figure~\ref{fig:tvod} presents willingness-to-pay for an
individual TVOD rental of a new Tamil film, prior to its
inclusion in any SVOD subscription.

\begin{figure}[h]
\centering
\includegraphics[width=0.85\textwidth]{fig_20_tvod.png}
\caption{Willingness to pay per TVOD rental of a new Tamil film.}
\label{fig:tvod}
\end{figure}

The modal willingness-to-pay band is INR 101--200 (34\%), followed by
INR 50--100 (31\%) and INR 201--300 (16\%).  Only 6\% of respondents
report willingness-to-pay above INR 300, indicating a soft ceiling
in the current market.  The aggregate weighted-mean willingness-to-pay
is approximately INR 142.

The finding has a direct supply-side implication.  The post-COVID
Tamil-language TVOD-rental market is \emph{viable} but priced
significantly below the comparable Hindi or English-language
TVOD-rental average reported by \citet{ficciey2024,kpmg2023}.

\section{Hypothesis Test Summary (H1--H5)}
\label{sec:hypothesis_summary_5}

Table~\ref{tab:hyp_5} summarises the results of the hypothesis tests
addressed in this chapter.

\begin{table}[h]
\centering
\caption{Hypothesis test summary -- viewer-side hypotheses}
\label{tab:hyp_5}
\small
\begin{tabular}{p{4.0cm} p{4.0cm} p{2.0cm} p{2.5cm} p{1.5cm}}
\toprule
\textbf{Hypothesis} & \textbf{Test} & \textbf{Stat} & \textbf{$p$} & \textbf{Status}\\
\midrule
H1: Theatre attendance fell      & Paired $t$-test & $t = 26.3$ & $<0.001$ & Supported \\
H2: OTT subscriptions rose       & Paired $t$-test; McNemar & $t = 24.8$; $\chi^2 = 76.5$ & $<0.001$ & Supported \\
H3: Content-driven choice rose   & McNemar & $\chi^2 = 49.7$ & $<0.001$ & Supported \\
H4: Genre distribution shifted   & Chi-square homogeneity & $\chi^2 = 89.4$ & $<0.001$ & Supported \\
H5: Demographic moderators       & Factorial ANOVA, regression & --- & --- & Partially supported (see Section~\ref{sec:moderators}) \\
\bottomrule
\end{tabular}
\end{table}

\section{Demographic Moderators of Adoption Magnitude (H5)}
\label{sec:moderators}

We tested the moderating role of age, location-tier, income and
education on the magnitude of post-COVID OTT adoption (defined as
$\Delta$ subscriptions held).  Multiple-regression results
(Table~\ref{tab:moderators}) show three statistically significant
predictors: \emph{age band} (younger respondents adopted more,
$\beta = -0.21$, $p < 0.001$), \emph{income band} (higher income
adopted more, $\beta = 0.18$, $p = 0.002$), and \emph{location-tier}
(Chennai and Tier-1 respondents adopted more,
$\beta = 0.13$, $p = 0.014$).  Education did not show a statistically
significant moderating effect once income was controlled
($\beta = 0.04$, $p = 0.41$).

\begin{table}[h]
\centering
\caption{Multiple regression -- predictors of OTT-adoption magnitude
(\(\Delta\) subscriptions 2019--2024)}
\label{tab:moderators}
\small
\begin{tabular}{lcccc}
\toprule
\textbf{Predictor}  & \textbf{$\beta$} & \textbf{SE} & \textbf{$t$} & \textbf{$p$}\\
\midrule
(Constant)            &  1.93 & 0.18 & 10.7 & $<0.001$ \\
Age band              & $-$0.21 & 0.05 & $-$4.2 & $<0.001$ \\
Income band           &  0.18 & 0.06 &  3.0 & 0.002 \\
Location-tier (Chennai/Tier-1) &  0.13 & 0.05 &  2.5 & 0.014 \\
Education             &  0.04 & 0.05 &  0.8 & 0.410 \\
Pre-COVID OTT subs    & $-$0.09 & 0.04 & $-$2.4 & 0.018 \\
\midrule
$R^2$ & \multicolumn{4}{l}{0.187, adjusted $R^2 = 0.176$, $F(5,394) = 18.1$, $p<0.001$}\\
\bottomrule
\end{tabular}
\end{table}

\textbf{Implication.}  H5 is partially supported.  Age, income and
location-tier moderate adoption magnitude as expected from
\citet{rogers2003} and \citet{shukla2022}; education does not, once
income is controlled.  The 18.7\% of variance explained by the
demographic block is modest, indicating that important non-demographic
variables (content-fit, social influence, network ties) drive
adoption but have not been measured here.  This gap is flagged for
future research in Section~\ref{sec:future} of
Chapter~\ref{ch:conclusion}.

\section{Mediation Analysis: OTT Adoption as Mechanism of
Theatre-Attendance Decline}
\label{sec:mediation}

A natural follow-up question is whether the decline in theatre
attendance (Section~\ref{sec:theatre_results}) is mechanistically
\emph{caused} by the rise in OTT adoption (Section~\ref{sec:ott_results}),
or whether the two are parallel effects of a common underlying
shock.  We address this with a Baron-Kenny-style mediation analysis
\citep{tashakkori2010,creswell2018} in which OTT-subscription growth
($\Delta$ subs) is treated as a candidate mediator between an
exogenous COVID-shock indicator (operationalised as ``did the
respondent's household experience a $\geq$ 4-month break in cinema
attendance during 2020--2021?'') and the change in monthly theatre
visits.

Step 1 (regress mediator on shock): the shock indicator significantly
predicts $\Delta$ subscriptions, $b = 0.84$, $p < 0.001$.
Step 2 (regress outcome on shock alone): the shock indicator
significantly predicts $\Delta$ theatre visits, $b = -1.19$,
$p < 0.001$.  Step 3 (regress outcome on shock and mediator
simultaneously): the direct effect of the shock falls to
$b' = -0.71$, $p < 0.001$, and the mediator is significant
$b = -0.57$, $p < 0.001$.  A bootstrapped Sobel test
\citep{johnson2007} on 5{,}000 resamples returns
$z = -7.4$, $p < 0.001$, with 95\% confidence interval
$[-0.65, -0.32]$, supporting partial mediation.  Approximately
40\% of the total effect of the shock on theatre attendance is
transmitted through OTT-adoption growth; the remaining 60\% is
direct (i.e., not channelled through OTT-subscription change but
through other unmeasured pathways such as habit erosion, content
disenchantment or price-perception shifts).

The mediation finding has two implications.  First, the
``substitution'' reading of OTT vis-\`a-vis theatre is partially
correct -- but only partially.  Second, the residual direct effect
indicates that even households which did not adopt OTT subscriptions
during the pandemic experienced theatre-attendance erosion, pointing
to non-OTT mechanisms (habit, content disenchantment, ticket-price
sensitivity) that warrant follow-up research
(Section~\ref{sec:future} of Chapter~\ref{ch:conclusion}).

A second mediation analysis on the genre-shift outcome similarly
finds that OTT-platform usage partially mediates the relationship
between the shock and the post-pandemic preference for thriller /
realistic / anthology genres, with approximately 56\% of the total
effect transmitted through platform-recommendation exposure
(operationalised as a self-report scale of ``how much of my Tamil-film
viewing is driven by what the platform recommends''), $z = 5.8$,
$p < 0.001$.  This is consistent with the platform-studies
literature on algorithmic curation \citep{nieborg2018,
gillespie2018,vandijck2018,plantin2018} and provides a partial
mechanism for the genre migration documented in
Section~\ref{sec:genre_results}.

\section{Joint Findings and Patterns}
\label{sec:joint_findings_5}

Several findings reinforce one another and warrant explicit synthesis:

\begin{itemize}[itemsep=2pt]
\item The decline in theatre attendance and the rise in OTT
      subscriptions are not in 1:1 substitution; OTT growth is
      larger in absolute viewing-hours-equivalent than the theatre
      decline, indicating \emph{net expansion} of total film-viewing
      time per household.
\item The shift from star-driven to content-driven choice is
      reinforced by the genre migration toward thriller / realistic /
      indie clusters, which by their nature are less star-dependent.
\item The shift toward solo viewing on smartphone aligns with the
      time-of-day shift toward late-night and the rise of the
      ``content-seeking'' gratification cluster
      \citep{sundet2021,lotz2022}.
\item The drop in self-reported piracy is consistent with the
      simultaneous rise in OTT subscriptions and willingness to
      pay; the affordable substitute exists in a form that
      previously did not.
\end{itemize}

\section{Chapter Summary}
\label{sec:viewer_summary}

This chapter has analysed the 400-respondent viewer-survey data
through 17 substantive variables, each compared between a 2019
baseline and a 2024 post-pandemic measure.  Hypotheses H1--H4 are
strongly supported; H5 is partially supported.  The post-pandemic
Tamil-film viewer attends fewer theatres, holds more OTT
subscriptions, watches more hours of OTT, prefers thriller /
realistic / anthology genres, makes content-driven (rather than
star-driven) choices, watches more solo and on smartphone or smart
TV, and pays a moderate amount for TVOD rentals when offered.  The
next chapter examines how the supply side of the industry has
reconfigured itself to serve this transformed viewer.
